\section*{الفصل \ref{ch:g10:inertia} --- القوى ومبدأ العطالة}

\begin{solution}{exo:g10:inertia:1}
(أ) \omterm{def:g10:universal-gravitation:weight}{الوزن} (عن بعد)، \omterm{def:g10:inertia:inventory}{ورد الفعل الناظمي} للطاولة (تماس).
(ب) \omterm{def:g10:universal-gravitation:weight}{الوزن} (عن بعد)، \omterm{def:g10:inertia:inventory}{وشدّ} السلك (تماس).
(ج) \omterm{def:g10:universal-gravitation:weight}{الوزن} (عن بعد)، \omterm{def:g10:inertia:inventory}{ورد الفعل الناظمي} للجليد (تماس)؛ \omterm{def:g10:inertia:inventory}{والاحتكاك}
مهمَل. (د) \omterm{def:g10:universal-gravitation:weight}{الوزن} وحده --- فلا شيء يلامس التفاحة.
\end{solution}

\begin{solution}{exo:g10:inertia:2}
الكتاب: $P = 0.250 \times 9.81 \approx \qty{2.45}{N}$، أي سهم طوله نحو
\qty{2.5}{cm}. والشخص: $P = 60 \times 9.81 \approx \qty{589}{N}$
--- أي سهم طوله \qty{5.89}{m}. فاختر مقياسًا يسع أكبر \omterm{def:g10:inertia:force}{قوة} على
الصفحة.
\end{solution}

\begin{solution}{exo:g10:inertia:3}
\omterm{def:g10:inertia:force}{نقطة التأثير} $A$؛ والاتجاه أفقي، متجه إلى اليسار؛
والشدة $3.5 \times 2 = \qty{7.0}{N}$.
\end{solution}

\begin{solution}{exo:g10:inertia:4}
\omterm{def:g10:universal-gravitation:weight}{الوزن} $\vect P$: عمودي، إلى أسفل، $P = 1.2 \times 9.81 \approx
\qty{11.8}{N}$. \omterm{def:g10:inertia:inventory}{ورد الفعل الناظمي} $\vect R$ للرفّ: المستقيم نفسه،
إلى أعلى، $R = \qty{11.8}{N}$ (فالقاموس ساكن، ومن ثم
تتعادل القوتان).
\end{solution}

\begin{solution}{exo:g10:inertia:5}
(أ) خطأ: \omterm{def:g10:inertia:inventory}{الاحتكاك} هو الذي يوقف الأشياء؛ فمن دون أي \omterm{def:g10:inertia:force}{قوة} يحفظ الجسم
حركته المستقيمة المنتظمة. (ب) صحيح (مبدأ العطالة).
(ج) خطأ: بل تؤثر فيه قوى \emph{\omterm{def:g10:inertia:compensate}{متعادلة}} (\omterm{def:g10:universal-gravitation:weight}{كالوزن}\ 
ورد الفعل مثلًا). (د) صحيح --- وهو القراءة النقيضة العكسية
للمبدأ.
\end{solution}

\begin{solution}{exo:g10:inertia:6}
\emph{1.} $v = 24/3.0 = \qty{8.0}{m/s}$. والقوى: \omterm{def:g10:universal-gravitation:weight}{الوزن} \omterm{def:g10:inertia:inventory}{ورد الفعل
الناظمي}؛ وهما يتعادلان (\omterm{def:g10:inertia:uniform}{حركة مستقيمة منتظمة}، \omterm{def:g10:inertia:inventory}{والاحتكاك}
مهمَل). \emph{2.} لم يعد \omterm{def:g10:inertia:inventory}{احتكاك} الإسمنت
مهمَلًا: فيكسب الجرد \omterm{def:g10:inertia:force}{قوة} خلفية غير متعادَلة،
وتتغير السرعة فعلًا --- فيتباطأ القرص ويتوقف.
\end{solution}

\begin{solution}{exo:g10:inertia:7}
\omterm{def:g10:universal-gravitation:weight}{الوزن} $P = 70 \times 9.81 \approx \qty{687}{N}$ إلى أسفل؛ \omterm{def:g10:inertia:force}{وقوة} الأرضية
إلى أعلى. \omterm{def:g10:inertia:uniform}{والحركة مستقيمة منتظمة}، فتتعادلان: أي أن الأرضية تدفع
\omterm{def:g10:inertia:force}{بقوة} \qty{687}{N}. ولا يتغير شيء في حالة السكون --- فالسكون والحركة المنتظمة
يعطيان \omterm{def:g10:inertia:diagram}{مخطط القوى} نفسه.
\end{solution}

\begin{solution}{exo:g10:inertia:8}
\emph{1.} في السكون، يتعادل \omterm{def:g10:inertia:inventory}{الشدّ} مع \omterm{def:g10:universal-gravitation:weight}{الوزن}:
$T = 1.8 \times 9.81 \approx \qty{17.7}{N}$. \emph{2.} ويحمل كل سلك
عمودي النصف: $T = \qty{8.8}{N}$.
\end{solution}

\begin{solution}{exo:g10:inertia:9}
عند الكبح: يحفظ جسم المسافر سرعته (بالعطالة) بينما تفقد
الحافلة سرعتها، فيتحرك المسافر إلى الأمام \emph{بالنسبة إلى
الحافلة}. وعند الانعطاف: يميل الجسم إلى مواصلة السير مستقيمًا بينما تنعطف الحافلة
تحته، فينجرف نحو خارج المنعطف. وفي
الحالتين تكون ``\omterm{def:g10:inertia:force}{القوة}'' وهمية؛ والناقص \omterm{def:g10:inertia:force}{قوة} حقيقية
(تشبّث، أو مقبض) تغيّر سرعة المسافر مع سرعة
الحافلة.
\end{solution}

\begin{solution}{exo:g10:inertia:10}
\omterm{def:g10:inertia:uniform}{الحركة مستقيمة منتظمة}، فيتعادل \omterm{def:g10:universal-gravitation:weight}{الوزن} ومقاومة الهواء: أي أن المقاومة
$= P = 85 \times 9.81 \approx \qty{834}{N}$، عمودية إلى أعلى.
\end{solution}

\begin{solution}{exo:g10:inertia:11}
\emph{1.} \omterm{def:g10:inertia:uniform}{الحركة مستقيمة منتظمة}: فتتعادل القوى الأفقية، ومن ثم تكون
المقاومة \qty{300}{N}، معاكسة للحبل. \emph{2.} \omterm{def:g10:universal-gravitation:weight}{وزن} المتزلج
ودفع الماء الصاعد على الزلاجتين.
\end{solution}

\begin{solution}{exo:g10:inertia:12}
\emph{1.} $P = 2.4 \times 9.81 \approx \qty{23.5}{N}$؛ وتحمل المركّبتان
العموديتان \omterm{def:g10:universal-gravitation:weight}{الوزن}: $2T\cos 40^\circ = P$، ومنه
$T = 23.5/(2 \times 0.766) \approx \qty{15.4}{N}$. \emph{2.}
يكبر $T = P/(2\cos\theta)$ بلا حدّ حين $\theta \to 90^\circ$
لأن $\cos\theta \to 0$؛ أما السلكان الأفقيان تمامًا فليست لهما مركّبة
عمودية \emph{البتة}، فلا يحمل الوزنَ شيء.
\end{solution}

\begin{solution}{exo:g10:inertia:13}
الدفع: \omterm{def:g10:universal-gravitation:weight}{الوزن}، ورد الفعل، ودفع اليد، \omterm{def:g10:inertia:inventory}{واحتكاك} صغير؛ وتزداد السرعة،
ومن ثم لا تتعادل القوى
(\cref{prop:g10:inertia:converse}). الانزلاق: \omterm{def:g10:universal-gravitation:weight}{الوزن}، ورد الفعل، \omterm{def:g10:inertia:inventory}{واحتكاك}
صغير؛ ويتباطأ الحجر، فلا تعادل مرة أخرى --- \omterm{def:g10:inertia:inventory}{والاحتكاك} هو
الفائض غير المتعادَل (وعلى جليد مثالي كان سينزلق أبدًا). السكون:
\omterm{def:g10:universal-gravitation:weight}{الوزن} ورد الفعل وحدهما، متعادلين
(\cref{thm:g10:inertia:principle}).
\end{solution}

\begin{solution}{exo:g10:inertia:14}
\emph{1.} بمبدأ العطالة، لا يقتضي حفظ سرعة ثابتة
أي \omterm{def:g10:inertia:force}{قوة} البتة --- ولا تكاد تؤثر \omterm{def:g10:inertia:force}{قوة} في الفضاء بين النجوم.
أما أرسطو فكان سيتنبأ بتوقف المسبار حين يكفّ
``المحرّك''. \emph{2.} السنة الواحدة $\approx \qty{3.156e7}{s}$، ومنه
$d = 17 \times \num{3.156e7} \approx \qty{5.4e8}{km} \approx
\qty{3.6}{au}$. \emph{3.} السرعة متجهة: فتغيير اتجاهها
يغيّر السرعة، وهو ما يقتضي \omterm{def:g10:inertia:force}{قوة} غير متعادَلة.
\end{solution}

\begin{solution}{exo:g10:inertia:15}
في السكون تتعادل القوى الثلاث:
$\vect F_3 = -(\vect F_1 + \vect F_2)$، ومنه $\vect F_3\,(-2, -4)$،
وشدتها $\sqrt{4 + 16} = \sqrt{20} \approx \qty{4.5}{N}$، متجهة
عكس محصلة القوتين الأوليين (إلى أسفل اليسار على الشبكة).
\end{solution}

\begin{solution}{pb:g10:inertia:1}
\textbf{1.} \omterm{def:g10:inertia:uniform}{حركة مستقيمة منتظمة} بسرعة \qty{0.75}{m/s}. والقوى: \omterm{def:g10:universal-gravitation:weight}{الوزن}
\omterm{def:g10:inertia:inventory}{ورد الفعل الناظمي} للسير (ولا حاجة إلى \omterm{def:g10:inertia:inventory}{احتكاك}: فالحقيبة
لا تنزلق على السير).
\textbf{2.} نعم --- \omterm{def:g10:inertia:uniform}{فالحركة مستقيمة منتظمة}
(\cref{thm:g10:inertia:principle}). و$P = 23 \times 9.81 \approx
\qty{226}{N}$، ومنه $R = \qty{226}{N}$.
\textbf{3.} المخططان متطابقان. فالسكون \omterm{def:g10:inertia:uniform}{والحركة المستقيمة
المنتظمة} حالة ميكانيكية واحدة بعينها: والعطالة لا
تفرّق بينهما.
\textbf{4.} تحفظ الحقيبة سرعتها \qty{0.75}{m/s} بينما يتوقف السير:
فتنزلق أو تنقلب إلى الأمام حتى يمتصّ \omterm{def:g10:inertia:inventory}{الاحتكاك} حركتها.
ولم يدفعها شيء --- إنما حفظت السرعة التي كانت لها.
\textbf{5.} تتجمع السرعتان في الاتجاه نفسه: $0.75 + 1.2 =
\qty{1.95}{m/s}$. وسرعتها ثابتة، ومن ثم تتعادل القوى المؤثرة فيها.
\textbf{6.} $P = 0.160 \times 9.81 \approx \qty{1.57}{N}$ إلى أسفل،
و$R = \qty{1.57}{N}$ إلى أعلى، و\emph{لا} \omterm{def:g10:inertia:force}{قوة} أفقية.
\textbf{7.} يحفظ القرص \qty{8.0}{m/s} دون أي محرّك على الإطلاق.
أما السطوح اليومية فتخفي \omterm{def:g10:inertia:inventory}{الاحتكاك} --- وهو \omterm{def:g10:inertia:force}{قوة} خلفية غير متعادَلة ---
فتبدو الحركة وكأنها تموت من تلقاء نفسها.
\textbf{8.} تتغير سرعته، فلم تعد \omterm{def:g10:inertia:compensate}{القوى متعادلة}
(\cref{prop:g10:inertia:converse})؛ والجاني هو \omterm{def:g10:inertia:inventory}{احتكاك}
الجليد الخشن.
\textbf{9.} لا: فالسرعة \emph{متجهة}، واتجاهها
يتغير، ومن ثم لا تتعادل القوى --- إذ يدفع الجليد جانبيًّا
على نصلَي حذائها.
\textbf{10.} (أ)--(3)، (ب)--(1)، (ج)--(2).
\textbf{11.} \omterm{def:g10:universal-gravitation:weight}{الوزن} $P = 80 \times 9.81 \approx \qty{785}{N}$، والمقاومة
معدومة: فلا شيء يعادل \omterm{def:g10:universal-gravitation:weight}{الوزن}، ومن ثم يجب أن تتغير السرعة ---
فيتسارع القافز نازلًا.
\textbf{12.} ما دامت المقاومة $<$ \omterm{def:g10:universal-gravitation:weight}{الوزن}، تبقى القوى غير \omterm{def:g10:inertia:compensate}{متعادلة}:
فتظل السرعة تنمو، وتنمو المقاومة معها.
\textbf{13.} عند السرعة الثابتة: يتعادل \omterm{def:g10:universal-gravitation:weight}{الوزن} والمقاومة، ومن ثم تكون المقاومة
\qty{785}{N}. \omterm{def:g10:inertia:uniform}{والحركة مستقيمة منتظمة} --- أي جسم ساقط في
الحالة الميكانيكية عينها التي للحقيبة على السير: عطالة
متنكّرة.
\textbf{14.} تتجاوز المقاومة الآن \omterm{def:g10:universal-gravitation:weight}{الوزن}: فلا تتعادل القوى، ومن ثم
تتغير السرعة --- فيتباطأ القافز، وهو لا يزال نازلًا؛ وكلما
هبطت السرعة تقلصت المقاومة عائدة نحو \omterm{def:g10:universal-gravitation:weight}{الوزن}.
\textbf{15.} وتثبت السرعة مجددًا: فالمقاومة $= P = \qty{785}{N}$ --- أي
المقدار نفسه بالضبط الذي كان عند \qty{50}{m/s}.
\textbf{16.} المقاومة الحدية تساوي \omterm{def:g10:universal-gravitation:weight}{الوزن} دائمًا، بمظلة
أو بغيرها. إنما تغيّر المظلة \emph{السرعة} التي عندها تبلغ المقاومة
\qty{785}{N}: أي \qty{5}{m/s} بدل \qty{50}{m/s}.
\textbf{17.} ``\dots القوى المؤثرة في الجسم \omterm{def:g10:inertia:compensate}{متعادلة} (أو لا تؤثر فيه
\omterm{def:g10:inertia:force}{قوة})'' --- وهما قراءتا \cref{thm:g10:inertia:principle}. والسكون
ليس إلا $\vect v = \vect 0$، أي \omterm{def:g10:inertia:uniform}{حركة مستقيمة منتظمة} كأي حركة أخرى.
\textbf{18.} الحلبة: فالقرص ينساب ولا محرّك في الأفق. ويشير غاليليو
إلى السطح المتزايد ملاسة --- \omterm{def:g10:inertia:inventory}{فالاحتكاك}، لا الطبيعة، هو ما
يوقف الأشياء.
\textbf{19.} بل تؤثر قوى --- \omterm{def:g10:universal-gravitation:weight}{الوزن} \qty{785}{N} إلى أسفل والمقاومة
\qty{785}{N} إلى أعلى؛ وهما \emph{متعادلتان}، وهذا بالضبط سبب كون
السرعة ثابتة.
\textbf{20.} السير: \omterm{def:g10:universal-gravitation:weight}{الوزن} \omterm{def:g10:inertia:inventory}{ورد الفعل الناظمي}. والحلبة: \omterm{def:g10:universal-gravitation:weight}{الوزن} \omterm{def:g10:inertia:inventory}{ورد
الفعل الناظمي}. والمظلة: \omterm{def:g10:universal-gravitation:weight}{الوزن} والمقاومة. والقانون الواحد: مبدأ
العطالة، بقراءتيه المباشرة والعكسية.
\end{solution}
